\begin{env}[4.1.4]
\label{0.4.1.4-fr}
Se dice que un haz de anillos $\mathcal{A}$ sobre un espacio topológico $X$
es \emph{reducido} en un punto $x$ de $X$ si el tallo $\mathcal{A}_x$ es un
anillo \emph{reducido} (1.1.1). Se dice que $\mathcal{A}$ es \emph{reducido}
si lo es en todo punto de $X$. Recordemos que un anillo $A$ se llama
\emph{regular} si cada uno de los anillos locales $A_{\mathfrak p}$ (donde
$\mathfrak p$ recorre el conjunto de los ideales primos de $A$) es un anillo
local regular. Diremos que un haz de anillos $\mathcal{A}$ sobre $X$
es \emph{regular en un punto} $x$ (resp. \emph{regular}) si el tallo
$\mathcal{A}_x$ es un anillo regular (resp. si $\mathcal{A}$ es regular
en todo punto). Por último, diremos que un haz de anillos $\mathcal{A}$ sobre
$X$ es \emph{normal en un punto} $x$ (resp. \emph{normal}) si su tallo
$\mathcal{A}_x$ es un anillo \emph{íntegro e integralmente cerrado} (resp. si
$\mathcal{A}$ es normal en todo punto). Diremos que un espacio anillado
$(X,\mathcal{A})$ tiene una de las propiedades anteriores si el haz de anillos
$\mathcal{A}$ tiene esa propiedad.

Un haz de \emph{anillos graduados} $\mathcal{A}$ es, por definición, un
haz de anillos que es suma directa (G, II, 2.7) de una familia
$(\mathcal{A}_n)_{n\in\mathbf Z}$ de haces de grupos abelianos, con las
condiciones $\mathcal{A}_m\mathcal{A}_n\subset\mathcal{A}_{m+n}$. Un
$\mathcal{A}$-\emph{Módulo graduado} es un $\mathcal{A}$-Módulo
$\mathcal{F}$ que es suma directa de una familia
$(\mathcal{F}_n)_{n\in\mathbf Z}$ de haces de grupos abelianos y satisface
las condiciones $\mathcal{A}_m\mathcal{F}_n\subset\mathcal{F}_{m+n}$. Equivalentemente,
se tiene
$(\mathcal{A}_m)_x(\mathcal{A}_n)_x\subset(\mathcal{A}_{m+n})_x$
(resp.
$(\mathcal{A}_m)_x(\mathcal{F}_n)_x\subset(\mathcal{F}_{m+n})_x$)
en todo punto $x$.
\end{env}
